% Owner: Member 3 (Nguyen Thanh Binh). Reviewers: Members 1 and 2 (Task.md, Report Ownership).
% The rules described here are owned by docs/ExperimentContract.md and implemented by tools/experiment/; this section reports them for readers of the paper.
% Text in \textbf{[...]} marks a value that exists only after the main campaign has run; replace it from results/<experiment-id>/ and never estimate it.

\section{Experimental setup}\label{sec:setup}

\subsection{Instances}\label{sec:setup-instances}

The experiments use 12 graphs from the Harwell--Boeing collection, listed in Table~\ref{tab:setup-instances}. % Member 2: add a \ref to the dataset provenance section here.
They are an initial subset of the 113 instances used in the reference paper, so the results do not support conclusions about the full collection. Every vertex is kept, including isolated vertices. Before each run, the input file is compared with the SHA-256 checksum recorded in the dataset manifest, and a mismatch stops the experiment.

\begin{table}[width=.9\linewidth,cols=5,pos=h]
\caption{Instances. $UB_0$ is the bandwidth of the identity labeling, which every method starts from. The SAT methods skip three instances, for the memory reason given in Section~\ref{sec:setup-conditions}.}\label{tab:setup-instances}
\begin{tabular*}{\tblwidth}{@{} LRRRL @{}}
\toprule
Instance & $n$ & $m$ & $UB_0$ & SAT runs \\
\midrule
ibm32 & 32 & 90 & 26 & yes \\
bcspwr01 & 39 & 46 & 38 & yes \\
bcsstk01 & 48 & 176 & 35 & yes \\
bcspwr02 & 49 & 59 & 34 & yes \\
curtis54 & 54 & 124 & 44 & yes \\
ash85 & 85 & 219 & 39 & yes \\
bcspwr03 & 118 & 179 & 115 & yes \\
bcsstk04 & 132 & 1758 & 47 & yes \\
bcspwr04 & 274 & 669 & 265 & yes \\
bcspwr05 & 443 & 590 & 435 & skipped \\
bcsstk20 & 485 & 1325 & 20 & skipped \\
sherman1 & 1000 & 1375 & 100 & skipped \\
\bottomrule
\end{tabular*}
\end{table}

\subsection{Methods and software}\label{sec:setup-methods}

Table~\ref{tab:setup-methods} lists the six compared configurations: four model or solving-scheme variants, two of which run on two backends each. Every backend runs in the environment pinned in the repository's \texttt{requirements.txt}, with Python~3.12.3.

\begin{table}[width=.9\linewidth,cols=3,pos=h]
\caption{Compared methods and the backend versions used.}\label{tab:setup-methods}
\begin{tabular*}{\tblwidth}{@{} L >{\raggedright\arraybackslash}p{.29\linewidth} >{\raggedright\arraybackslash}p{.33\linewidth} @{}}
\toprule
Method & Model and solving scheme & Backend \\
\midrule
\texttt{sat\_order\_repeated} & Order encoding, new solver for each bound & CaDiCaL 1.5.3~\cite{cadical} through PySAT 1.9.dev22~\cite{pysat} \\
\texttt{sat\_order\_incremental} & Order encoding, one solver for all bounds & The same CaDiCaL build \\
\texttt{cp\_ortools} & CP formulation & OR-Tools CP-SAT 9.15.6755~\cite{ortools} \\
\texttt{cp\_cpo} & The same CP formulation & IBM CP Optimizer, CPLEX Optimization Studio 22.2~\cite{cplex} \\
\texttt{ilp\_cplex} & ILP formulation (Section~\ref{sec:ilp}) & IBM CPLEX 22.2.0.0~\cite{cplex} \\
\texttt{ilp\_gurobi} & The same ILP formulation & Gurobi 14.0.0~\cite{gurobi} \\
\bottomrule
\end{tabular*}
\end{table}

\subsection{Machine and run conditions}\label{sec:setup-conditions}

All runs execute on one laptop with an Intel Core i5-13450HX processor (10 physical and 16 logical cores) and 15{,}698~MiB of RAM, under Ubuntu 24.04.5 LTS with Linux 6.8.0. Only one run executes at a time, and every backend uses one thread.

Each method runs three times on each instance, with repetition $r \in \{1, 2, 3\}$ using seed $r - 1$ on every backend that accepts a seed. CaDiCaL accepts no seed through PySAT, so its seed is recorded as unavailable. Instances run in manifest order and repetitions in ascending order, and the method order of repetition $r$ is the order of Table~\ref{tab:setup-methods} rotated left by $r - 1$, so no method always runs first.

Every run has a budget of 300~seconds for model construction and optimization together. The timer starts after the input is loaded and its checksum verified, and it excludes process start-up, the final independent check, and writing the results. All SAT bound tests share the one budget, and each CP and ILP backend receives the time remaining after model construction. A run that has not ended 60~seconds after its deadline is killed; its last saved incumbent is kept and the kill is recorded as its termination reason. Each run executes in a new process, and peak memory is the largest resident set size of that process or of any process it started, including the memory of the native solver libraries.

The SAT order encoding has $n \cdot n(n-1)/2$ clauses for the all-different constraint, so its memory grows with $n^3$. On bcspwr04 (274 vertices, 10.3 million clauses), the SAT search reached a peak of 5.4--5.7~GiB within 60~seconds. Scaled by clause count, bcspwr05, bcsstk20 and sherman1 would need about 23, 31 and 270~GiB, more than the machine's RAM, so the SAT methods skip them. These skips are not runs and not timeouts. In total, the campaign has 198 runs (12 instances for the CP and ILP methods and 9 for the SAT methods, three repetitions each), with at most 16.5 hours of budgeted solving time.

\subsection{Result status and checking}\label{sec:setup-status}

Every labeling that a method returns or saves during its search is checked by one shared checker, independent of all solver models. The checker confirms that the labeling is a permutation of $\{1, \ldots, n\}$ and recomputes its bandwidth, which must equal the bandwidth the method claims. Every method starts from the identity labeling $f_0$, with upper bound $UB_0$ and the common lower bound $LB$ (1 for a graph with an edge). Each run ends with one of four statuses:
\begin{itemize}
\item \texttt{OPTIMAL}: a checked labeling together with a proof that no smaller bandwidth exists. For SAT, the proof is an UNSAT answer at the bound one below the incumbent, or an incumbent equal to $LB$. For CP and ILP, it is the backend's optimality result with a certified bound equal to the incumbent.
\item \texttt{FEASIBLE}: a checked labeling without such a proof, usually because the time limit was reached.
\item \texttt{UNKNOWN}: no checked labeling and no proof, for example after an interruption before initialization.
\item \texttt{ERROR}: invalid output, a failed check, an exception, a crash, or a solver result that contradicts itself.
\end{itemize}
Timeouts and interruptions are never treated as UNSAT. Before the performance runs, the CP and ILP methods returned the optimum on paths, cycles, complete graphs, disconnected graphs and edgeless graphs of at most 8 vertices, compared with an exhaustive search over all labelings. \textbf{[SAT: the same check is not yet automated; state its result once Member 1 adds it.]}

\subsection{Reproducibility}\label{sec:setup-reproducibility}

The campaign ran at commit \textbf{[campaign commit]} of the project repository on \textbf{[campaign dates]}. For every run, the repository stores the configuration, the environment record, the exact dataset manifest, one JSON record per run, the checked labeling, and the solver log, under \texttt{results/\textbf{[experiment id]}/}. The whole campaign can be repeated with the command in the repository's \texttt{AGENTS.md}; the time limits make absolute times depend on the machine.